\section{The Three Readings Compared}
\label{sec:comparison}

The three readings hold less in common than their shared Gospel suggests, and
what they share can be said exactly. All three read the Postcommunion's
\latin{oboedire mandatis} as the answer to the Introit's \latin{non
oboedivimus}: in the first the word received becomes the word obeyed, in the
second obedience is the exiles' road home, and in the third the law given to
Israel is asked of all who receive the gifts. All three begin from the Introit's
confession, though each hears it differently, as the table below shows; and all
three give the Communion's hope in God's word a place: in the first it is the
prayer of one who hopes in a promised word, as the father did; in the second,
the Church's comfort in her humiliation; in the third, the nations' hope in the
word first given to Abraham.

Beyond that the common ground is partial. The first and the third take the
Gospel's word, \latin{Filius tuus vivit}, as the formulary's hinge, and both
rest on the Fathers' agreement that the healing at Capharnaum was done by a
word, without the bodily presence the father asked for: the first as the object
of one man's faith, the third as a figure of the two peoples' relation to that
word. The second is not built on the Gospel. It centres on the Introit, the
Offertory and the Epistle, and gives the Gospel a supporting place, as one
exile's crisis in a world that St Gregory says now sends us to God. The first
and the second hear the Introit's confession in the order St Augustine finds in
the saints, God's justice praised first and his gifts asked after; the third
hears it instead as Israel's confession, prayed by the Church as her own. The
witnesses overlap in the same uneven way. St Augustine speaks in all three, for
the order of prayer, the prepared heart and the physician's due time in the
first, for the rivers of Babylon and the road of the law in the second, for the
figure of the two peoples in the third; St Jerome's innocent youths, who
confess \latin{ex persona populi}, stand in the second and third; and St Thomas
Aquinas gives the first its stages of faith and the third its interceding
Fathers.

They differ in the question they put to the texts, and so in which elements
weigh most in each.

\begin{comparisontable}{Question}{The word that heals}{The song of the exiles}{The two peoples}
The question it answers & What faith does Christ ask of the father, and why does
he heal by a word without going down? & How is a people that has sinned and
lives in exile to pray, sing and spend its time? & Whose story is the healing at
Capharnaum, and whose voices does the Mass let us hear?\\
Where it centres & The Gospel's exchange: the plea to come down, the rebuke, the
word, the going & The Introit, the Offertory and the Epistle & The Gospel as a
figure, with the chants as voices\\
Whose situation & One believer's, and every believer's & The Church's and each
soul's present exile & The history of two peoples\\
What the Introit is & The confession that rightly opens petition & The exile's
confession of God's justice, with the road home in its verse & Israel's
confession, prayed by the Church as her own\\
What the Offertory is & Memory of what is not seen: faith in absence & The
Church's lament and her memory of Sion & Israel's lament by Babylon's rivers\\
The strongest difficulty & The hardness of the rebuke to a dying boy's father
& The Introit's ``because we have sinned'', and the close of Ps~136 & That the
figure can be heard as setting Israel against the Church\\
Carried by & Gregory and Chrysostom & Augustine and Hilary & Augustine and
Aquinas\\
\end{comparisontable}

No doctrine of one reading is denied by another, and the differences inside the
readings are real: whether the ruler believed before the word (Augustine says
not; Gregory says he believed and doubted), why Christ would go to the
centurion and not to the ruler (Gregory, to humble the pride that honours rank;
Chrysostom, because the centurion's faith was already perfect), whether the
willows of Babylon are barren men or saints, whether the little ones of Ps~136
are sins at their birth or persons brought to Christ, and whether the ruler
figures the people that needed signs or the patriarch who prays for it. Between
the three readings the relation is one of level and scope. Aquinas himself
holds the first and the third together, the literal and moral growth of one
man's faith and the allegory of the Fathers and their people, at different
senses of the same text. The second makes no claim about the other two, but
read beside them it gives them a setting: the father who walks home on a word
and the nations who believe without signs are both on the road of exiles who
remember Sion.

The readings meet one real tension in the chants. Fromage's two choirs give the
Introit and Offertory to Israel; Augustine's sermon on Ps~136 gives the
Offertory to the Church in her own captivity. The third reading hears Israel's
voice first and the second hears the Church's, and Jerome's youths, who confess
for a people they belong to, show how one voice can carry both: the Church
prays Israel's confession as her own because she prays it, as Azarias did, for
and with a people.

Each reading accounts for something the others leave aside. The first alone
takes the Gospel's dialogue phrase by phrase, and it alone reads the Secret's
medicine and the Gradual's due time as the father's situation carried into the
Mass. The second alone gives full weight to the two verbs of singing that the
Epistle commands and the Alleluia performs, and it alone faces the end of the
Offertory's psalm. The third alone hears the
Gospel's place in its chapter, between the believing Samaritans and the
Galileans who wanted signs, and alone gives the nations' universality in the
Gradual and the Alleluia's unsung ``among the nations'' a voice.

The Gospel's last words, \latin{et credidit ipse, et domus eius tota}, show
where the first and third meet and where the second stands apart. In the first
they are the end of one man's growth in faith, from need, through trust in a
word, to belief with his house. In the third they are a figure of the whole
house, the people whose Fathers prayed for it and the nations grafted beside it,
for whom St Cyril hopes the end of the story: ``then shall all Israel be
saved''. The second does not take up those words. Its use of the Gospel goes
no further than the father's journey: the world full of death and grief, which
St Gregory says now sends us to God, sends him to Christ, and the word that
reaches his son is the divine word St Jerome finds in the furnace. That
reading's answer is not a household's faith but the exiles' song, sung to God
in a strange land by a heart that waits for the city it remembers.
